\chapter{Operating the local causal service and qualifying its limits}
\label{in:operations}

\section{An execution identity is larger than a request name}

A durable analysis must bind the canonical request, scientific context, input
bytes, code, schema, numerical environment, seed and budget. A repeated job
name alone cannot establish that a result is still valid. The implemented
runtime records these identities and rejects idempotency collisions involving
changed requests or execution identities.~\citep{v12_in_runtime}

Artifacts use canonical JSON, content and descriptor hashes, opaque identities
and bounded local reads. Regular-file and nonsymlink checks reduce accidental
path redirection. A content hash establishes byte identity under the relevant
assumption about the hash function; it is not an authenticated endorsement of
the input's scientific accuracy. A local operator can still supply a false
declaration whose bytes remain perfectly stable.

The scientific request contains no arbitrary shell command, Python import or
network endpoint. Fixed worker entry points execute the supported operations.
This makes the scope of the local tool explicit and allows a review grant to
refer to actual operations. The current controller's trusted caller object is
not a substitute for an institutional identity and approval service.

\section{Budgets, durable jobs and cancellation}

Scientific workers have wall-time and address-space limits. The memory bound
applies to the worker, not to all memory used by the resident SDK and
coordinator. A complete operating measurement must include both. Small control
or error envelopes may exceed an extremely small requested output limit to
explain why the request cannot be fulfilled. These are disclosed boundaries of
the present implementation.~\citep{v12_in_runtime}

Report jobs retain a hashed prefix of completed work identities. The local
retry rule is bounded to three attempts under the original deadline. Publication
is atomic, and replay checks request, artifact and checkpoint identities before
continuation. The recorded restart fixture kills a report worker after its
first checkpoint and completes all one hundred distinct work identities once.
This is evidence for that local workflow, not a distributed exactly-once
guarantee for arbitrary external side effects.

Cancelling an in-flight protocol request and cancelling a durable application
job are different operations. The inspected specification also recognizes
cancellation/completion races.~\citep{v12_mcp_legacy_cancellation} If publication
has already completed, a cancellation response must report that terminal state.
It must not imply that a displayed result or an external action was undone.
The present service has no operation that changes a bank account or sends a
customer communication.

\section{Fault evidence and what remains unmeasured}

Nine recorded fault cases cover replay with conflicting budgets, wall timeout,
output limits, worker memory, kill/restart with checkpoint, cooperative
cancellation, changed artifact bytes, changed scientific context and changed
checkpoints.~\citep{v12_in_runtime} The tests distinguish preserving a result from
merely restarting a process. An interrupted computation that silently uses new
inputs is a different analysis even if it eventually finishes.

These fixtures do not measure institution-scale availability, multi-host
recovery, adversarial tenant isolation or a service-level commitment. They also
do not show that streaming always improves elapsed time. The earlier complete
local workload comparisons retain small or absent gains; engineering choices
must be justified by the intended workload rather than by the existence of a
streaming interface.

Operational qualification should measure the complete request: artifact import,
scientific work, uncertainty, serialization, reporting and recovery. A benchmark
of one reducer omits fitting and coordination costs. The comparison needs
matched inputs, numerical tolerance and an explicit accounting of unsuccessful
requests, just as a scientific comparison needs a common target.

\section{Compatibility is a maintained scientific obligation}

A new backend release can change defaults, graph encodings, estimator targets,
uncertainty or failure behavior. A compatible JSON schema alone does not prove
scientific compatibility. The proposed adapter register should bind a version,
accepted inputs, method identity, preserved result types and matched adverse
fixtures. A material semantic change requires a reviewed successor rather than
relabeling an old qualification.

Transport migration likewise needs real client/server tests. The current legacy
MCP choice is useful for its tested environment. Supporting a modern protocol
or both eras requires a new compatibility argument, discovery behavior and
failure tests. It should not be advertised by extending a version string while
keeping the old handshake.

Ownership must also remain explicit. The causal library owns its result
contract and adapters. Numerical providers own their implementations and
scientific domains. Consumers own their target interpretation and displayed
decision. The bank owns actual data authority, assignment records, operational
constraints and acceptance. A report can connect those responsibilities; it
cannot replace their evidence.

\section{What a recovery identity must preserve}

Consider a daily deposit report that aggregates customer-level evidence into
a 90-day funding decision. A restarted process must answer the same question
using the same admitted observations. Replacing a data extract with a later
vintage, changing a sign convention, or upgrading a numerical routine can
alter the scientific result even when the file name and command remain
unchanged. This is why recovery needs an analysis identity rather than only
a process identifier.

Write the intended identity as
\[
 J=\operatorname{Hash}(Q,A,S,E,\rho,B),
\]
where $Q$ is the scientific request, $A$ the input identities, $S$ the schema
and method semantics, $E$ the relevant execution environment, $\rho$ the
randomization state where applicable, and $B$ the declared budget. This is a
proposed completeness condition. The inspected runtime fingerprints a selected
list of source files and schemas, Python, and four named package versions.
Its request digest omits the display request identifier and idempotency key
while retaining the other request fields. It rejects reuse of a key with a
changed request or selected execution identity
\cite{exec13_ops_source}.

The selected fingerprint is useful but does not describe every numerical or
platform dependency. It does not, for example, independently record the
complete transitive dependency graph, numerical-library build, operating
system and hardware state. A future numerical adapter must identify which
of those properties can affect its result, then bind and test them. Extending
the fingerprint is an engineering obligation; it does not establish that the
underlying causal assumptions are true. Conversely, irrelevant display
metadata need not invalidate a deterministic result if the scientific and
execution identities remain unchanged.

Three compatibility questions follow. Byte replay asks whether the same
admitted computation returns the same canonical bytes. Numerical compatibility
allows a declared tolerance or stochastic comparison under an unchanged
mathematical target. Scientific compatibility asks whether the population,
treatment, horizon, units, assumptions and inferential meaning remain the
same. A backend upgrade can satisfy a numerical tolerance while changing
which observations are dropped or which population is weighted. Such an
upgrade needs a semantic review even if its output schema is unchanged.

For the funding report, sensible adverse fixtures include a changed date
cutoff, a currency-unit conversion, a negative monetary map, a missing
assignment field, an unbounded uncertainty set and a default that drops
customers without complete histories. The comparison must preserve valid
refusals and ambiguous results. Substituting a point estimate for an
unbounded set is not a compatibility repair.

\section{Checkpoint invariants and the scope of idempotence}

Suppose a report combines deterministic contributions $s_i$ for a frozen
ordered list of work identities $(a_1,\ldots,a_m)$. A useful checkpoint is
\[
 C_k=(J,(a_1,\ldots,a_k),(s_1,\ldots,s_k),h_k),
\]
where $h_k$ binds the preceding contents. The required invariant is that the
completed identifiers form a prefix of the frozen request, each identifier
has one committed contribution, and the public result is a function of those
committed contributions. Recomputing an uncommitted contribution after a
crash may waste work; committing it twice changes the answer.

An exact project construction makes the distinction visible. Three intended
contributions are $100,-40,70$, so the total is $130$. A crash and retry can
produce the calculation trace $100,-40,-40,70$. Naively appending every
returned value gives $90$. A map from work identity to its immutable
contribution still gives $130$, provided a repeated identity cannot carry
changed contents. These are finite arithmetic results, not measurements of
a distributed service~\citep{exec13_ops_exact}.

The inspected local report worker verifies its completed prefix and digest,
skips already committed identities, and updates summaries and completed
identities together in the job state. Its preserved restart trial completed
100 distinct work identities, with 100 unique entries, after two worker
starts. The more recent preserved fault replay also reports the nine declared
cases passing; that replay was not a fresh MCP client/server trial
\cite{exec13_ops_source,exec13_ops_faults}. This evidence supports the named
local report path. It does not establish recovery for arbitrary estimators,
multi-host execution or transactions at a bank.

The atomicity boundary is also specific. The artifact store publishes its
content and metadata as individually atomic files, checks any already
present bytes, and can complete a repeated deterministic publication. The
report's result reference and terminal job state are then published in the
job record. These are several transitions, not one transaction spanning
arbitrary external systems. A crash can leave an unreferenced artifact even
when no incorrect public result is exposed. Retention and garbage collection
therefore need a separate rule that preserves artifacts referenced by admitted
jobs and evidence. Hash verification alone cannot establish that cleanup is
safe.

A customer communication illustrates the stronger problem. If an external
system sends an offer and the acknowledgement is lost, a retrying client may
not know whether the offer was sent. Recording the local calculation once
does not resolve that ambiguity. An eventual external-action adapter would
need a cooperating transaction or deduplication contract, reconciliation and
a named owner for unresolved outcomes. The present service performs no such
bank action. Its proposal should retain report generation as a distinct,
read-oriented capability rather than infer action safety from a successful
restart.

\section{Retries purchase attempts, not independent chances of success}

Let $K$ be a maximum number of attempts and $q$ the probability of a transient
failure on each attempt. If failures are independent and identically
distributed, the probability of obtaining a result is $1-q^K$ and the
expected number of attempts is
\[
 \mathbb E N_K=\sum_{j=0}^{K-1}q^j.
\]
For the stipulated values $q=1/5$ and $K=3$, these are $124/125=0.992$
and $31/25=1.24$. The numbers express an assumed failure law. Nine passing
fault fixtures do not estimate $q$, and a retry cap alone cannot establish
the 99.2\% success probability~\citep{exec13_ops_exact}.

Persistent failures change the calculation. Suppose one tenth of requests
cannot succeed under the supplied environment, while the remaining nine
tenths have independent per-attempt failure probability one tenth. Three
attempts then succeed with probability
\[
 \frac9{10}\left(1-\frac1{10^3}\right)=0.8991,
\]
and require $0.1(3)+0.9(1+0.1+0.01)=1.299$ attempts on average.
The marginal first-attempt failure probability is $0.19$. Treating that
marginal value as independent across retries would predict
$1-0.19^3=0.993141$, substantially overstating success. The retry has not
changed an unavailable dependency, an unsupported method or contradictory
input declarations. Operational classification should distinguish those
persistent causes from a recoverable interruption.

A bounded retry policy also retains its original time account. In the
inspected report path, remaining time is computed from the job creation
timestamp and original wall budget, and the worker-start count is capped at
three. Restarting does not intentionally grant three new full deadlines.
The implementation uses local wall-clock timestamps and operating-system
timers; this is not a distributed clock or a hard real-time guarantee.
Launch overhead, host suspension, clock adjustment and recovery of a failed
machine require a target-environment study before making a service-level
claim~\citep{exec13_ops_source}.

Statistical retries need an additional distinction. Replaying an interrupted
calculation with preserved random state is different from drawing new
randomness until a desired conclusion appears. If three independent valid
null $p$-values are generated and the smallest is selected, the probability
of at least one value below $0.05$ is
\[
 1-(1-0.05)^3=0.142625.
\]
This exact counterexample assumes independent uniform null $p$-values;
it does not allege that the runtime performs such selection. It shows why
the proposed adapter must separate mechanical recovery from a new statistical
replicate and preserve all requested analyses. The same issue arises when
timeouts depend on difficult data and only easy completed fits enter a
coverage report~\citep{exec13_ops_exact}.

\section{Resource budgets and cancellation at the right boundary}

The local worker sets an address-space limit and a wall-time timer. The
coordinator starts a fixed worker entry point, collects its output and
validates the canonical response against the requested output size. This
supports a bounded set of local operations. It does not impose one aggregate
memory bound on the worker, coordinator, SDK and all clients, and the output
check is applied after collection rather than as an incremental bound on an
arbitrary producer's stream~\citep{exec13_ops_source}.

For a proposed third-party numerical adapter, admission should therefore
include peak resident memory of the whole workflow, temporary disk use,
serialization and parent-process buffers. A limit that is safe for a short
report may be inappropriate for an estimator returning a large covariance
matrix. If a result exceeds the public response size, the design can publish
a separately bounded artifact with explicit custody and expose a small
reference. It must retain the scientific object's meaning and cannot silently
truncate interval components or omit unsuccessful fits to meet a size limit.

Cancellation is a race between a request and the computation's next state
transition. For pure report generation, a useful rule is that an accepted
cancellation recorded before the publication transition prevents a new public
result; cancellation after
publication returns the actual completed state. The inspected implementation
checks a cancellation flag around report work and before publication. That
cooperative behavior does not guarantee that every numerical kernel can be
interrupted at an arbitrary instruction. A later adapter may need worker
termination and a rule for discarding uncommitted artifacts.

A timeout should be recorded as an execution outcome. A scientifically
unidentified effect can instead be a completed result with an explicit
scientific limitation. Converting both into one generic failure statistic
obscures the remedy: the first may need resource or implementation work, while
the second may need a different target, new information or abstention. A
successful process exit likewise says nothing by itself about whether the
effect is precise enough to support a decision.

\section{Ownership and a reversible adoption path}

Operational ownership follows the claim being made. A method maintainer owns
the supported algorithm and numerical domain. The common-library owner owns
the request/result contract, method-version binding and admitted conversions.
The consumer owns the displayed target, units and business action. A later
bank study additionally needs a data owner for the actual observation and
assignment processes and a decision owner for losses and feasible actions.
These are proposed responsibilities, not names or approvals that have already
been supplied.

For example, a deposit report may use a valid backend estimate while a
consumer interprets the coefficient as currency rather than thousands of
currency units. The remedy belongs partly in the shared unit contract and
partly in the consuming application's acceptance examples. A central
numerical team cannot certify a silently altered presentation. Conversely,
an application owner cannot declare a confidence procedure valid merely
because the resulting interval is displayed correctly.

The proposal should make an adoption decision reversible at its natural
boundaries. Retain the incumbent expert workflow and its artifacts while a
new adapter is evaluated on matched public/generated requests. Publish both
methods' complete outputs with their own identities; compare discrepancies
before replacing the default. A schema migration should have a declared
mapping and refusal cases for information it cannot preserve. Where a new
method changes the estimand, the migration is a new scientific analysis,
not a formatting conversion. A rollback retains access to the previously
admitted version and evidence without rewriting its result as if produced
by the new version.

The same approach applies to workspace consumers. Read-only inspection of a
repository can reveal candidate objects and interfaces. It cannot establish
that its owner accepts the proposed contract, that its launch environment is
compatible, or that an integration saves work. A local mock consumer can test
the proposed request and response boundary now. Actual consumer execution,
deployment ownership and institutional operation remain distinct later
claims, with their own narrowly defined studies.

\section{The economic denominator is every attempted request}

Let $R$ denote whether a request returns a usable result and $D$ its incremental
decision value relative to the incumbent workflow, including the specified
fallback when $R=0$. Let $C$ include all request-specific compute, review,
retry and recovery costs. The per-request economic target is
\[
 \mathbb E[D-C],
\]
plus a consistent allocation of fixed construction and shared operating
costs. Reporting $\mathbb E[D\mid R=1]$ with a nominal successful-run cost
answers a different question. Reliability and economic comparisons require
both the success probability and what happens on the unsuccessful branch.

In an exact stipulated example, 70\% of requests are easy, complete and add
USD~10 of gross value. Another 15\% complete but add USD~$-30$. The remaining
15\% fail and revert to the baseline action, adding zero gross value but
requiring USD~20 of recovery work. Every attempt costs USD~2 in tools and
computation. The completed-only calculation reports
\[
 \frac{0.70(10)+0.15(-30)}{0.85}-2
 =\frac{16}{17}\simeq0.941
\]
USD per completed request. The all-attempt account instead gives
\[
 0.70(10)+0.15(-30)-2-0.15(20)=-2.5
\]
USD per attempted request. The apparent gain changes sign when the
unsuccessful work and correct denominator enter the account. These numbers
are generated planning assumptions, not workflow prices or measured bank
returns~\citep{exec13_ops_exact}.

Measurement should therefore retain attempted task identity, terminal
execution state, scientific disposition, usable-result status, fallback
action and disjoint cost categories. Repeated attempts are costs within a
request, not new independent banking opportunities. Correlated requests from
one institution, calendar or scientific task family also cannot be treated
as independent evidence merely because their request identifiers differ.
The analysis unit should follow the source of independent information.

A prototype can measure machine resources and trace whether required fields
survive recovery. It cannot price saved analyst time without observing a
matched workflow, nor infer business value from seconds saved. The next
inexpensive operational study should compare complete local workflows under
declared faults and budgets, preserve adverse results, and report source and
environment identities. Only a subsequent justified platform or bank study
can supply institution-specific availability, ownership, labor, and
decision-value evidence. Their absence limits those claims while leaving
the public/generated proposal work executable.

\section{The next complete qualification}

The remaining engineering route is concrete: choose source-audited packages
and permitted integration forms, implement representative bridges, and run
matched consumer workflows across all three surfaces. The trial must include
unsupported methods, late observations, graph ambiguity, weak or disconnected
effect sets, interruption and changed-input rejection. Passing results should
name the actual supported environments and domains.

This work will close an engineering gap only in those domains. Negative
persistent-map, latent-graph and dependent-learner studies remain scientific
obligations. Independent review, authentic bank validation and measured
economics remain additional gates. A coherent runtime makes their required
inputs and outputs reviewable; it does not turn missing information into an
accepted causal answer.

\label{exec13:operations-end}
